\section{Formal results}
\label{app:formal}

This appendix details the formal results summarized in
\cref{sec:design:abstract:formal}.
\Cref{thm:design:embed,thm:design:findend}
are mechanized in Lean~4 (\cref{app:lean,app:boundary});
\cref{thm:design:verify} has a paper proof only.

\subsection{Always-embeddability}
\label{app:formal:embed}

Two closure properties of $L$ do the work:

\begin{lemma}[Concatenation closure]
\label{lem:design:concat}
If $a \in L$ and $b \in L$, then $a || b \in L$.
\end{lemma}

\begin{lemma}[Wrapping]
\label{lem:design:wrap}
If $m \in L$ and $(o,c) \in \Pi$, then $o||m||c \in L$.
\end{lemma}

Concatenation follows by induction on the derivation of $L$,
and wrapping directly from the nesting rule (details in \cref{app:lean}).

To capture the shape of the delimiter strings $s_o,s_c$
that practical hosts use,
we define a \emph{matched delimiter pair} inductively:

\begin{definition}[Matched delimiter pair]
\label{def:design:delim}
The pair $(s_o,s_c)$ is \emph{matched} if it can be formed by these rules:
\begin{itemize}
\item	$(o,c)$ is matched for every $(o,c) \in \Pi$;
\item	if $(s_o,s_c)$ is matched and $(o,c) \in \Pi$,
	then $(o||s_o,\; s_c||c)$ is matched; and
\item	if $(s_o,s_c)$ is matched and $x \in N$ is a nonmatcher,
	then both $(x||s_o,\;s_c)$ and $(s_o,\;s_c||x)$ are matched.
\end{itemize}
\end{definition}

Intuitively, $s_o$ supplies one or more open matchers
and $s_c$ the corresponding close matchers in properly-nested order,
with nonmatcher characters interspersed freely except next to $m$,
exactly the shape of practical delimiters such as
the \texttt{M} and \texttt{"} surrounding the parenthesis
in \verb|M"(|$\dots$\verb|)"|.

\begin{proof}[Proof sketch of \cref{thm:design:embed}]
By induction on the derivation that $(s_o,s_c)$ is matched
(\cref{def:design:delim}).
The base case wraps $m$ using \cref{lem:design:wrap},
and a matcher step likewise wraps the induction hypothesis;
a nonmatcher step prepends or appends a single-character
nonmatcher string (which is in $L$ by the flat rule)
using \cref{lem:design:concat}.
\end{proof}

\Cref{thm:design:embed} establishes the compositional core
of escapeless embedding:
for \emph{any} valid matchertext $m$, with no exceptions,
the delimited result $s_o||m||s_c$ is \emph{itself} valid matchertext.
It can therefore be embedded again, to any depth,
with no escaping and no expansion of the inner string,
the property that distinguishes matchertext
from the quoting and escaping mechanisms surveyed in \cref{sec:bg}.
The complementary claim of \cref{sec:design:abstract:embed},
that a host parser scanning left-to-right can \emph{unambiguously}
locate the terminating $s_c$,
follows from the same balanced structure;
\cref{app:formal:findend} makes it operational
and \cref{thm:design:findend} mechanizes it.

\subsection{Deciding matchertext}
\label{app:formal:verify}

Whereas \cref{thm:design:embed} shows that embedding \emph{produces}
valid matchertext,
a host language, or a checkpoint guarding against untrusted input
(\cref{sec:bg}),
must also be able to \emph{recognize} it.
Membership in the language $L$ of \cref{sec:design:abstract}
is decided by a single linear scan,
shown as \cref{alg:design:verify}.
The stack holds the openers seen so far that are not yet matched;
each closer must pop the opener of its own pair,
and the string is matchertext exactly when
the scan finishes with an empty stack.

\begin{algorithm}[!t]
\caption{Deciding whether a string $s$ is matchertext.}
\label{alg:design:verify}
\begin{algorithmic}[1]
\Function{Verify}{$s$}
	\State $st \gets$ empty stack \Comment{unmatched openers}
	\For{each character $x$ in $s$}
		\If{$x \in O$} push $x$ onto $st$
		\ElsIf{$x \in C$}
			\State \textbf{if} $st$ is empty \textbf{or} $(\textsc{Top}(st), x) \notin \Pi$
				\textbf{then return} \textsc{false}
			\State pop $st$
		\EndIf
	\EndFor
	\State \Return whether $st$ is empty
\EndFunction
\end{algorithmic}
\end{algorithm}

\begin{theorem}[Correctness of \textsc{Verify}]
\label{thm:design:verify}
For every string $s$,
\cref{alg:design:verify} returns \textsc{true} if and only if $s \in L$.
\end{theorem}

\begin{proof}[Proof sketch]
By induction on the scan, the stack holds exactly the currently unmatched openers,
so \textsc{Verify} rejects exactly the unmatched or mismatched matchers.
\end{proof}

The procedure runs in $O(n)$ time and $O(d)$ space
on an input of length $n$ with maximum nesting depth $d$,
which \cref{sec:eval:nesting} finds is almost always tiny.
This is the concrete basis for the security argument of \cref{sec:bg}:
admitting untrusted embedded content requires only this passive linear check,
rather than a content-modifying transformation
that must anticipate every sequence sensitive to the host language.
It is also exactly the check that our empirical analysis (\cref{sec:eval})
applies to every string literal and comment in the corpus.

\subsection{Locating the end of an embedding}
\label{app:formal:findend}

A host must also find where an embedding $s_o||m||s_c$ ends
without being misled by matchers inside $m$.
A host language that has already validated $m$ as matchertext,
for example because $m$ was produced by
the process in \cref{thm:design:embed},
can scan the text after $s_o$ with a depth counter starting at $0$
and stop at the first closer that makes the depth negative.
By \cref{thm:design:findend} this is the first character of $s_c$,
which \cref{def:design:delim} makes a closer.
A host language reading untrusted input, by contrast,
must instead use the typed stack of
\cref{alg:design:verify},
since a depth counter alone accepts wrong-kind pairs such as \verb|([)]|.
The host must then check that the complete $s_c$ follows.
